\chapter{Etat de l'art et positionnement du projet}

\guided{Ce chapitre doit d\'emontrer une revue critique de la litt\'erature ou des solutions existantes, et non une simple compilation. Appuyez-vous prioritairement sur des articles index\'es, conf\'erences, normes, brevets, documentations techniques de r\'ef\'erence. Terminez par le positionnement explicite du projet.}

\section{Contexte scientifique et technique}
\todo{Introduire le domaine et les enjeux du sujet.}

\section{Analyse critique des travaux existants}
\todo{Comparer 3 \`a 5 approches, outils ou architectures pertinentes.}

\begin{table}[H]
\centering
\caption{Exemple de tableau comparatif de l'\'etat de l'art}
\begin{tabularx}{\textwidth}{>{\RaggedRight\arraybackslash}p{2.8cm} >{\RaggedRight\arraybackslash}X >{\RaggedRight\arraybackslash}p{3cm} >{\RaggedRight\arraybackslash}p{2.8cm}}
\toprule
R\'ef\'erence / solution & Principe / apport & Limites identifi\'ees & Int\'er\^et pour le projet \\
\midrule
Solution A & \todo{...} & \todo{...} & \todo{...} \\
Solution B & \todo{...} & \todo{...} & \todo{...} \\
Solution C & \todo{...} & \todo{...} & \todo{...} \\
\bottomrule
\end{tabularx}
\end{table}

\section{Positionnement du projet}
\todo{Montrer le besoin non couvert, la contribution attendue et la justification du choix retenu.}

\guided{Attendus forts pour une lecture type ASIIN/EUR-ACE : pertinence des sources, qualit\'e du recul critique, lien clair entre lacunes identifi\'ees et choix du projet.}
